\documentclass{beamer}
\usetheme{CambridgeUS}
\usecolortheme{beaver}

\usepackage[utf8]{inputenc}
\usepackage[spanish]{babel}
\usepackage{amsmath}
\usepackage{amssymb}
\usepackage{graphicx}
\usepackage{booktabs}

\title{Framework de Explicabilidad Agnóstica al Modelo}
\subtitle{Para la Predicción de Series Temporales Financieras}
\author{Rushell Vanessa Zavalaga Orozco}
\institute{Escuela de Ciencia de la Computación\\
Universidad Nacional de San Agustín de Arequipa}
\date{}

\begin{document}

% -----------------------------------------------------------------------------
\begin{frame}
    \titlepage
\end{frame}

% -----------------------------------------------------------------------------
\begin{frame}{Motivación (1/2)}
    \small
    \begin{itemize}
        \item Los modelos de aprendizaje profundo predicen con alta
        precisión series temporales financieras.
        \vspace{0.15cm}
        \item Pero \textbf{no revelan} qué información del historial
        determinó cada predicción.
        \vspace{0.15cm}
        \item En entornos regulados, la transparencia no es opcional:
        \begin{itemize}
            \item SR~11-7 (Reserva Federal) exige entender capacidades
            y límites del modelo antes de usarlo
            \item El Reglamento de IA de la UE exige transparencia
            interpretable en usos financieros críticos
            \item Sin explicabilidad no es posible detectar sesgos
            ni fallos del modelo
        \end{itemize}
        \vspace{0.15cm}
        \item Un modelo que no puede explicarse no puede auditarse:
        la opacidad impide distinguir si aprendió una relación
        económica real o una coincidencia espuria en los datos de
        entrenamiento.
    \end{itemize}
\end{frame}

% -----------------------------------------------------------------------------
\begin{frame}{Motivación (2/2)}
    \small
    \begin{itemize}
        \item Un modelo entrenado sobre datos financieros históricos
        hereda los sesgos de ese historial.
        \vspace{0.2cm}
        \item Por eso este trabajo reporta también los resultados que
        \textbf{no favorecen a la propuesta} (persistencia ingenua,
        referencia por hoja), en vez de omitirlos.
        \vspace{0.2cm}
        \item Al no requerir gradientes, además, reduce la barrera de
        cómputo para auditar modelos, un punto de acceso equitativo
        para instituciones con menos recursos.
        \vspace{0.4cm}
        \item \textbf{Pregunta central:} ¿qué variable fue relevante
        para una predicción y \textit{cuándo} en el historial lo fue?
    \end{itemize}
\end{frame}

% -----------------------------------------------------------------------------
\begin{frame}{El Problema}
    \small
    \begin{block}{Problema de cómputo}
        Un mecanismo de explicación agnóstico al modelo (sin
        gradientes) y de costo fijo por instancia (sin optimizar una
        máscara para cada predicción).
    \end{block}
    \vspace{0.15cm}
    \begin{alertblock}{Por qué sigue sin resolverse}
        DynaMask~\cite{crabbe2021} optimiza una máscara por gradiente
        (hasta 1000 épocas, 5 hiperparámetros), excluye modelos no
        diferenciables como Random Forest. Sus extensiones directas,
        ExtrMask~\cite{enguehard2023} (ICML 2023) y
        ContraLSP~\cite{liu2024contralsp} (ICLR 2024), mejoran la
        perturbación pero mantienen la dependencia del gradiente.
    \end{alertblock}
    \vspace{0.15cm}
    \begin{center}
        \footnotesize Cinco años después, ningún método de esta línea
        funciona sin gradientes ni se ha validado en finanzas.
    \end{center}
\end{frame}

% -----------------------------------------------------------------------------
\begin{frame}{Trabajos Relacionados: Dos Ejes}
    \textbf{Eje 1 — XAI con dimensión temporal (sin agnosticismo total)}
    \begin{itemize}
        \small
        \item Huang et al.\ (2025) — ShapeX: Shapley sobre shapelets;
        solo clasificación.
        \item Yadav y Subbian (2025) — diagnostican por qué gradiente,
        oclusión y permutación fallan en predicción dinámica.
    \end{itemize}
    \vspace{0.2cm}
    \textbf{Eje 2 — Frameworks temporalmente conscientes y agnósticos}
    \begin{itemize}
        \small
        \item Tonekaboni et al.\ (2020) — FIT; Leung et al.\ (2023) —
        WinIT: pensados para predicción en línea (clínico), no ventana
        única.
        \item Crabbé y Van Der Schaar (2021) — DynaMask: salida
        $T{\times}V$, pero \textbf{requiere gradientes}.
        \item Bento et al.\ (2021) — TimeSHAP: solo modelos secuenciales.
        \item Raykar et al.\ (2023) — TsSHAP: features predefinidas,
        sin resolución por celda.
        \item Roshinta et al.\ (2026) y Kim y Park (2026, GS-SHAP):
        agnósticos, pero agregan variables (PCA / grupos) y evalúan un
        único tipo de modelo.
    \end{itemize}
\end{frame}

% -----------------------------------------------------------------------------
\begin{frame}{La Brecha}
    \begin{alertblock}{Ningún trabajo previo satisface simultáneamente}
        \begin{enumerate}
            \item Agnosticismo real al modelo (sin gradientes, sin
            estructura secuencial obligatoria)
            \item Resolución completa por \textbf{celda}
            $(t,v)$, sin agrupar variables ni segmentos
            \item Verificación empírica de que las explicaciones
            \emph{cambian} entre modelos de distinta naturaleza
            computacional (Adebayo et al., 2018)
        \end{enumerate}
    \end{alertblock}
    \vspace{0.3cm}
    \begin{block}{Hipótesis}
        Es posible construir un método de perturbación temporalmente
        ordenado, agnóstico al modelo, que produzca explicaciones
        consistentes con el comportamiento de cualquier arquitectura.
    \end{block}
    \begin{center}
        \textbf{Salida central}: matriz
        $\hat{\alpha} \in \mathbb{R}^{T \times V}$
    \end{center}
\end{frame}

% -----------------------------------------------------------------------------
\begin{frame}{Pregunta de Investigación e Hipótesis}
    \begin{block}{Pregunta de investigación}
        ¿Es posible atribuir la importancia de cada variable y cada
        paso de tiempo dentro de una ventana de entrada, de forma
        agnóstica al modelo, capturando el comportamiento real de
        arquitecturas de distinta naturaleza computacional?
    \end{block}
    \vspace{0.3cm}
    \begin{alertblock}{Hipótesis}
        Es posible construir un método de perturbación temporalmente
        ordenado, agnóstico al modelo, que produzca explicaciones
        consistentes con el comportamiento de cualquier arquitectura.
    \end{alertblock}
\end{frame}

% -----------------------------------------------------------------------------
\begin{frame}{Objetivo General}
    \begin{center}
        \vspace{0.5cm}
        \large
        Proponer un framework de explicabilidad agnóstico al modelo
        que atribuya importancia por variable y paso de tiempo
        ($\hat{\alpha} \in \mathbb{R}^{T \times V}$) en la predicción
        de series temporales financieras, sin depender de la
        estructura interna del modelo explicado.
    \end{center}
\end{frame}

% -----------------------------------------------------------------------------
\begin{frame}{Objetivos Específicos}
    \small
    \begin{enumerate}
        \item Diseñar una interfaz uniforme de predicción, mediante
        un patrón adaptador de dos niveles, que desacople el motor de
        perturbaciones de la estructura interna de cualquier modelo.
        \vspace{0.2cm}
        \item Formular un motor de perturbaciones temporales que
        atribuya importancia a nivel de celda $(t,v)$, preservando el
        orden y la coherencia temporal del resto de la ventana.
        \vspace{0.2cm}
        \item Entrenar modelos de distinta naturaleza computacional
        (LSTM, Transformer, Random Forest) sobre el S\&P 500 como
        banco de pruebas del framework.
        \vspace{0.2cm}
        \item Evaluar el framework frente al estado del arte
        agnóstico, en precisión de identificación, consistencia
        entre modelos y velocidad de cómputo.
    \end{enumerate}
\end{frame}

% -----------------------------------------------------------------------------
\begin{frame}{Contribución: Interfaz Uniforme de Predicción}
    \begin{block}{Patrón adaptador de dos niveles}
        El \texttt{TemporalPerturbationEngine} nunca accede al modelo
        directamente; solo conoce una función
        $f: \mathbb{R}^{N\times T\times V}\to\mathbb{R}^N$.
    \end{block}
    \begin{itemize}
        \item \textbf{Nivel 1 — Adaptación de firma}: cada modelo
        implementa un \emph{wrapper} propio.
        \begin{itemize}
            \small
            \item Random Forest: aplana $(T,V)\to(T{\cdot}V)$
            \item LSTM / Transformer: tensor PyTorch, inferencia sin
            gradientes
        \end{itemize}
        \item \textbf{Nivel 2 — Inyección de $f$}: el
        \texttt{ModelAgnosticExplainer} construye $f$ y la inyecta en
        el motor; cualquier modelo externo puede conectarse sin tocar
        el framework.
    \end{itemize}
    \vspace{0.2cm}
    \begin{center}
        \textit{Inversión de dependencias}: el framework no depende
        de los modelos, los modelos se adaptan al framework.
    \end{center}
\end{frame}

% -----------------------------------------------------------------------------
\begin{frame}{Datos y Modelos}
    \begin{columns}
        \begin{column}{0.48\textwidth}
            \begin{block}{Datos: S\&P 500}
                \begin{itemize}
                    \item Período: 2010--2024
                    \item 3{,}754 días hábiles
                    \item $V = 15$ variables:\\
                    5 base + 10 indicadores técnicos
                    \item Ventanas: $T=20$, $H=1$
                    \item Split temporal: 80\% desarrollo /
                    20\% prueba
                \end{itemize}
            \end{block}
        \end{column}
        \begin{column}{0.48\textwidth}
            \begin{block}{Modelos heterogéneos}
                \begin{itemize}
                    \item \textbf{LSTM}\\
                    2 capas, hidden=128, dropout 0.2
                    \item \textbf{Transformer}\\
                    2 capas, 4 heads,\\
                    attention pooling aprendido
                    \item \textbf{Random Forest}\\
                    200 árboles, sklearn\\
                    aplana $(T,V) \to (T{\times}V)$
                \end{itemize}
            \end{block}
        \end{column}
    \end{columns}
    \vspace{0.3cm}
    \begin{center}
        Interfaz uniforme: $(N, T, V) \rightarrow (N,)$ —
        compatible con PyTorch y scikit-learn.\\
        \small MSE prueba (target Close): LSTM 0.00342 \quad
        Transformer 0.00284 \quad RF 0.12131
    \end{center}
\end{frame}

% -----------------------------------------------------------------------------
\begin{frame}{Motor de Perturbaciones Temporales (1/2)}
    \begin{block}{Idea central}
        Para cada posición $(t, v)$ de la ventana de entrada,
        se reemplaza el valor por la referencia $r_v$ (media de
        entrenamiento) y se mide el cambio en la predicción:
    \end{block}

    \begin{equation*}
        \alpha(t, v) = \left| f(x) - f(x_{-(t,v)}) \right|
    \end{equation*}

    \begin{alertblock}{¿Por qué oclusión?}
        \small
        No requiere gradientes (funciona en Random Forest) y
        perturba una sola celda a la vez, preservando la coherencia
        temporal del resto de la ventana.
    \end{alertblock}
\end{frame}

% -----------------------------------------------------------------------------
\begin{frame}{Motor de Perturbaciones Temporales (2/2)}
    \begin{block}{Optimización: batch único}
        En lugar de $T \times V$ inferencias individuales, se
        construye un batch de $TV+1 = 301$ ventanas evaluado en una
        sola llamada al modelo. El costo escala $O(T\cdot V)$, no
        de forma independiente por dimensión.
    \end{block}

    \begin{block}{Normalización L1}
        \begin{equation*}
            \hat{\alpha}(t, v) = \frac{\alpha(t, v)}
            {\sum_{t'}\sum_{v'} \alpha(t', v')}
        \end{equation*}
    \end{block}

    \begin{center}
        \small Cada celda de $\hat\alpha$ indica la proporción del
        cambio total en la predicción atribuible a esa variable en
        ese paso de tiempo.
    \end{center}
\end{frame}

% -----------------------------------------------------------------------------
\begin{frame}{Pipeline del Framework}
    \begin{center}
        \includegraphics[width=0.95\linewidth]{PIPELINE.png}
    \end{center}
    \begin{center}
        \scriptsize
        (1) Datos $\to$ (2) Preparación / ventanas $\to$ (3) Modelo
        $\to$ (4) Interfaz de predicción $\to$ (5) Motor de
        perturbaciones $\to$ (6) Matriz $\hat{\alpha}$
    \end{center}
\end{frame}

% -----------------------------------------------------------------------------
\begin{frame}{Validación con Ground Truth Sintético}
    \begin{block}{¿Qué se hizo?}
        Objetivo de caja blanca $f(\mathbf{X}) = \sum_{(t,v)\in
        \mathcal{A}} x_{t,v}^2$ con celdas relevantes conocidas
        ($T{=}20$, $V{=}15$, $n{=}10$ repeticiones, 2 escenarios:
        \textit{Rare Feature} y \textit{Rare Time}).
    \end{block}
    \vspace{0.1cm}
    \centering
    \scriptsize
    \begin{tabular}{lcccc}
        \toprule
        \textbf{Método} & \textbf{AUP (RF)} & \textbf{AUR (RF)}
        & \textbf{AUP (RT)} & \textbf{AUR (RT)}\\
        \midrule
        FO (propuesto) & 0.9960 & 0.16 & 0.9960 & 0.16 \\
        FP             & 0.9960 & 0.18 & 0.9960 & 0.18 \\
        IG             & 0.9960 & 0.16 & 0.9960 & 0.16 \\
        SVS            & 0.9960 & 0.17 & 0.9960 & 0.15 \\
        \bottomrule
    \end{tabular}
    \vspace{0.2cm}
    \begin{alertblock}{Resultado}
        FO iguala a SVS/IG/FP y \textbf{supera} el AUP reportado por
        DynaMask (0.97) — sin requerir gradientes y ejecutándose
        también sobre Random Forest.
    \end{alertblock}
\end{frame}

% -----------------------------------------------------------------------------
\begin{frame}{Verificación de Agnosticidad sobre S\&P 500}
    \begin{columns}
        \begin{column}{0.48\textwidth}
            \begin{block}{Consistencia inter-modelo}
                Spearman $\rho$ entre matrices $\hat{\alpha}$,
                $N=100$ instancias.
                \vspace{0.2cm}
                \begin{tabular}{lc}
                    \toprule
                    \textbf{Par} & \textbf{$\rho$} \\
                    \midrule
                    LSTM -- Transf. & 0.748 \\
                    LSTM -- RF      & 0.255 \\
                    Transf. -- RF   & 0.382 \\
                    \bottomrule
                \end{tabular}
            \end{block}
        \end{column}
        \begin{column}{0.48\textwidth}
            \begin{block}{Significancia estadística}
                Mann-Whitney U vs.\ baseline aleatorio.
                \vspace{0.2cm}
                \begin{tabular}{lc}
                    \toprule
                    \textbf{Modelo} & \textbf{$p$} \\
                    \midrule
                    LSTM        & $<10^{-300}$ \\
                    Transformer & $<10^{-300}$ \\
                    Rand. Forest & $<10^{-300}$ \\
                    \bottomrule
                \end{tabular}
            \end{block}
        \end{column}
    \end{columns}
    \vspace{0.3cm}
    \begin{center}
        Neuronal--neuronal $>$ neuronal--ensamble: las explicaciones
        capturan el comportamiento de cada modelo, no un artefacto
        del método.
    \end{center}
\end{frame}

% -----------------------------------------------------------------------------
\begin{frame}{Resultados: Importancia por Variable}
    \begin{columns}
        \begin{column}{0.55\textwidth}
            \includegraphics[width=\linewidth]{results/variable_importance.png}
        \end{column}
        \begin{column}{0.42\textwidth}
            \begin{itemize}
                \item \textbf{LSTM y Transformer}: priorizan
                SMA\_5, SMA\_10, SMA\_20 y High.
                \vspace{0.2cm}
                \item \textbf{Random Forest}: concentra
                importancia en Close y Low.
                \vspace{0.2cm}
                \item Indicadores de oscilación (MACD, RSI,
                BB\_width) $\approx$ 0 en los tres modelos.
                \vspace{0.2cm}
                \item[] {\scriptsize \textit{Interpretar con cautela,
                ver diapositiva de retorno logarítmico.}}
            \end{itemize}
        \end{column}
    \end{columns}
\end{frame}

% -----------------------------------------------------------------------------
\begin{frame}{Resultados: Patrón Temporal de Importancia}
    \begin{columns}
        \begin{column}{0.55\textwidth}
            \includegraphics[width=\linewidth]{results/temporal_importance.png}
        \end{column}
        \begin{column}{0.42\textwidth}
            \begin{itemize}
                \item \textbf{LSTM}: decreciente pasado$\to$presente
                (razón 15.7) — memoria distribuida.
                \vspace{0.2cm}
                \item \textbf{Transformer}: perfil casi uniforme
                (razón 1.5) — attention pooling.
                \vspace{0.2cm}
                \item \textbf{Random Forest}: concentrado en lag~0
                (razón 860) — sin modelado secuencial.
            \end{itemize}
        \end{column}
    \end{columns}
\end{frame}

% -----------------------------------------------------------------------------
\begin{frame}{Resultados: Matrices de Importancia $\hat{\alpha}$}
    \begin{center}
        \includegraphics[width=0.9\linewidth]{results/heatmaps_comparison.png}
    \end{center}
    \begin{center}
        \small
        LSTM y Transformer: filas SMA concentran valores altos,
        distribuidos en el tiempo. Random Forest: importancia
        concentrada en lag~0 (Close, Low).
    \end{center}
\end{frame}

% -----------------------------------------------------------------------------
\begin{frame}{Robustez: Target de Retorno Logarítmico}
    \begin{alertblock}{¿Por qué?}
        El target Close tiene alta autocorrelación con el precio
        actual. Se repite todo el pipeline usando el retorno
        logarítmico como target, que rompe esa autocorrelación.
    \end{alertblock}
    \centering
    \scriptsize
    \begin{tabular}{lccc}
        \toprule
        \textbf{Modelo} & \textbf{MSE} & \textbf{MAE} & \textbf{RMSE}\\
        \midrule
        Persistencia (retorno=0) & 0.00257 & 0.03728 & 0.05067 \\
        LSTM         & 0.00268 & 0.03842 & 0.05178 \\
        Transformer  & 0.00260 & 0.03772 & 0.05104 \\
        Random Forest& 0.00446 & 0.05625 & 0.06677 \\
        \bottomrule
    \end{tabular}
    \vspace{0.2cm}

    \textbf{Ningún modelo supera la persistencia ingenua} —
    consistente con eficiencia de mercado débil. Aun así, el patrón
    de consistencia se reproduce: LSTM--Transf. $\rho{=}0.541$ $>$
    LSTM--RF $\rho{=}0.187$, Transf.--RF $\rho{=}0.311$
    ($p\approx 0$ vs.\ baseline aleatorio en los tres modelos).
\end{frame}

% -----------------------------------------------------------------------------
\begin{frame}{Robustez Adicional}
    \begin{itemize}
        \item \textbf{Exactitud direccional}: 46--50\% en los tres
        modelos y ambos targets — indistinguible del azar (50\%).
        \vspace{0.2cm}
        \item \textbf{Sensibilidad al vector de referencia}
        (media vs.\ mediana vs.\ cero): LSTM y Transformer estables
        ($\rho>0.88$); Random Forest más sensible ($\rho$ hasta
        0.666). El patrón neuronal $>$ ensamble se mantiene en las
        tres referencias.
        \vspace{0.2cm}
        \item \textbf{Régimen de mercado} (bajista 2022 vs.\ alcista
        2023--2024): correlaciones algo menores en el régimen bajista,
        pero el orden relativo (neuronal--neuronal $>$
        neuronal--ensamble) se preserva en ambos.
    \end{itemize}
    \begin{center}
        \includegraphics[width=0.75\linewidth]{results/consistency_robustness.png}
    \end{center}
\end{frame}

% -----------------------------------------------------------------------------
\begin{frame}{Consideración Específica para Random Forest}
    \begin{block}{Hipótesis probada}
        Random Forest particiona su entrada en hojas. En vez de una
        referencia externa (media/mediana/cero), usar una referencia
        \textbf{condicionada a la hoja} por instancia (promedio de
        ejemplos de entrenamiento co-hoja, ponderado por cercanía),
        motivado por Hooker et al.~\cite{hooker2021}: sin reentrenar.
    \end{block}
    \vspace{0.2cm}
    \begin{alertblock}{Resultado: negativo}
        \small
        Sintético: AUP/AUR sin cambio significativo (0.9376 vs.\
        0.9374). Real (S\&P 500): la consistencia con LSTM/Transformer
        \textbf{empeora} de $\rho{=}0.184/0.311$ a
        $\rho{=}{-}0.111/{-}0.232$.
    \end{alertblock}
    \vspace{0.2cm}
    \begin{center}
        La referencia condicionada a la hoja es \emph{interna} al
        modelo (rompe el criterio externo compartido con LSTM/
        Transformer). El framework mantiene una referencia externa
        idéntica en su naturaleza para los tres modelos: esta
        alternativa empeora la comparabilidad en vez de mejorarla.
    \end{center}
\end{frame}

% -----------------------------------------------------------------------------
\begin{frame}{Ablación: ¿Importa Respetar el Orden Temporal?}
    \begin{block}{Variante \emph{shuffled} (control negativo)}
        Sustituye la celda $(t,v)$ por el valor de la misma variable
        en \textbf{otro instante} $t'\neq t$ elegido al azar, en vez
        de la referencia fija $r_v$ — viola el orden temporal.
    \end{block}
    \centering
    \scriptsize
    \begin{tabular}{lcc}
        \toprule
        \textbf{Par de modelos} & \textbf{FO} & \textbf{Shuffled} \\
        \midrule
        LSTM -- Transformer          & 0.748 & 0.487 \\
        LSTM -- Random Forest        & 0.255 & $-0.055$ \\
        Transformer -- Random Forest & 0.382 & 0.251 \\
        \bottomrule
    \end{tabular}
    \vspace{0.2cm}
    \begin{alertblock}{Resultado}
        El benchmark sintético (determinista) \textbf{no distingue}
        FO de \emph{shuffled}. Solo sobre datos reales se observa el
        daño: consistencia cae, fidelidad $\Delta_\alpha$ se reduce
        hasta ${\times}38$ en RF, y las explicaciones dejan de ser
        reproducibles ($\rho\approx0.55$ entre corridas).
    \end{alertblock}
\end{frame}

% -----------------------------------------------------------------------------
\begin{frame}{Comparación en Condiciones Idénticas: FO vs.\ SVS}
    \begin{block}{¿Qué se hizo?}
        Se ocluyen las celdas más importantes según cada método y se
        mide el desplazamiento de predicción $\Delta_\alpha$
        (fidelidad). SVS = Shapley Value Sampling, único otro método
        agnóstico con salida $T{\times}V$ aplicable a los 3 modelos.
    \end{block}
    \begin{columns}
        \begin{column}{0.55\textwidth}
            \includegraphics[width=\linewidth]{results/exp3_mae_shift.png}
        \end{column}
        \begin{column}{0.42\textwidth}
            \scriptsize
            \begin{tabular}{lcc}
                \toprule
                \textbf{Modelo} & \textbf{FO} & \textbf{SVS} \\
                \midrule
                LSTM        & \textbf{0.5033} & 0.4891 \\
                Transformer & \textbf{0.4518} & 0.4352 \\
                Rand. Forest& 0.3047 & \textbf{0.3212} \\
                \bottomrule
            \end{tabular}
            \vspace{0.3cm}
            FO más fiel en modelos neuronales; SVS mejor en RF
            (interacciones combinatorias entre árboles).
        \end{column}
    \end{columns}
\end{frame}

% -----------------------------------------------------------------------------
\begin{frame}{Velocidad: FO vs.\ Estado del Arte}
    \centering
    \scriptsize
    \begin{tabular}{lccc}
        \toprule
        \textbf{Modelo} & \textbf{FO (ms)} & \textbf{SVS (ms)}
        & \textbf{Aceleración} \\
        \midrule
        LSTM          & 17.6  & 469.1  & ${\times}26.6$ \\
        Transformer   & 13.8  & 413.5  & ${\times}30.0$ \\
        Random Forest & 112.0 & 2519.4 & ${\times}22.5$ \\
        \bottomrule
    \end{tabular}
    \vspace{0.3cm}

    \begin{itemize}
        \item TimeSHAP~\cite{bento2021}: 2{,}813 ms/muestra$^\dagger$
        — FO es ${\times}159{,}8$ más rápido.
        \item GS-SHAP~\cite{gsshap2026}: 1{,}194 ms/muestra$^\dagger$
        — FO es ${\times}67{,}8$ más rápido.
    \end{itemize}
    \vspace{0.2cm}
    \begin{center}
        \includegraphics[width=0.7\linewidth]{results/timing_comparison.png}
    \end{center}
    \tiny $^\dagger$Cifras de Kim y Park (2026), Tabla C3 — preprint
    sin arbitrar; comparación provisional.
\end{frame}

% -----------------------------------------------------------------------------
\begin{frame}{Resumen de Avances frente al Estado del Arte}
    \centering
    \tiny
    \setlength{\tabcolsep}{3pt}
    \begin{tabular}{p{1.9cm}p{2.7cm}p{2.6cm}p{2.6cm}}
        \toprule
        \textbf{Trabajo previo} & \textbf{Brecha} & \textbf{Su resultado}
        & \textbf{Este trabajo} \\
        \midrule
        DynaMask~\cite{crabbe2021} & Gradientes; sin RF
        & AUP 0.97, sin RF & AUP \textbf{0.9960}, RF \checkmark \\
        TimeSHAP~\cite{bento2021} & Solo secuencial; lento
        & 2{,}813 ms/m.$^\dagger$ & \textbf{17.6} ms/m.\ (${\times}159.8$) \\
        GS-SHAP~\cite{gsshap2026} & Salida agrupada $K{\times}J$
        & 1{,}194 ms/m.$^\dagger$; solo $K{\times}J$
        & \textbf{17.6} ms/m.; $T{\times}V$ \checkmark \\
        \bottomrule
    \end{tabular}
    \vspace{0.3cm}
    \begin{alertblock}{En síntesis}
        Mismo o mejor AUP, resolución completa por celda, funciona
        sobre modelos no diferenciables, y órdenes de magnitud más
        rápido que las alternativas basadas en Shapley.
    \end{alertblock}
\end{frame}

% -----------------------------------------------------------------------------
\begin{frame}{Conclusiones}
    \begin{itemize}
        \item Framework de explicabilidad temporal agnóstico al
        modelo mediante un adaptador de dos niveles: cualquier modelo
        con firma $(N,T,V)\to(N,)$ se conecta sin modificar el motor.
        \vspace{0.2cm}
        \item Sobre datos sintéticos: FO iguala a SVS en precisión de
        identificación (AUP $=0.9960$).
        \vspace{0.2cm}
        \item Sobre S\&P 500: FO supera a SVS en fidelidad en modelos
        neuronales, ${\times}22$--${\times}30$ más rápido; SVS mejor
        en Random Forest.
        \vspace{0.2cm}
        \item Atribuciones estadísticamente no triviales
        ($p<10^{-300}$) y consistentes con la familia arquitectónica
        ($\rho=0.748$ neuronal--neuronal vs.\ $\rho<0.40$
        neuronal--ensamble), robustas a target, referencia y régimen
        de mercado.
    \end{itemize}
\end{frame}

% -----------------------------------------------------------------------------
\begin{frame}{Limitaciones}
    \begin{itemize}
        \small
        \item \textbf{Autocorrelación del target Close}: mitigada
        evaluando también el retorno logarítmico (Log\_Return).
        \item \textbf{Escalado MinMax} ajustado sobre 2010--2021 no
        cubre los máximos históricos de 2022--2024 — perjudica en
        particular a Random Forest.
        \item \textbf{Benchmark sintético no discriminativo}: función
        determinista sin ruido; con ruido aditivo, FO y SVS sí se
        diferencian (SVS más robusto por promediar permutaciones).
        \item \textbf{Alcance}: un único índice (S\&P 500). Evaluado
        frente a referencia, régimen de mercado y target alternativo,
        pero no frente a otro mercado.
    \end{itemize}
\end{frame}

% -----------------------------------------------------------------------------
\begin{frame}{Trabajo Futuro}
    \begin{itemize}
        \item Perturbación bidireccional (arriba/abajo de $r_v$) para
        recuperar el \textbf{signo} de la atribución, no solo la
        magnitud.
        \item Replicar el pipeline sobre otro índice (p.\ ej.\ NASDAQ
        Composite, \texttt{\^{}IXIC}) cambiando un solo parámetro de
        configuración.
        \item Correlacionar volatilidad realizada por ventana con la
        concordancia inter-modelo de $\hat{\alpha}$, instancia por
        instancia (no solo por régimen agregado).
        \item Extender a perturbación por bloques de celdas para
        escalar a $T$ o $V$ mucho mayores.
    \end{itemize}
\end{frame}

% -----------------------------------------------------------------------------
\begin{frame}{Declaración de Uso de IA}
    \begin{block}{Alcance del uso de asistentes de IA generativa}
        \begin{itemize}
            \item Apoyo en edición de redacción y verificación de
            referencias bibliográficas.
            \item Generación de código auxiliar para experimentos
            adicionales de robustez, siempre bajo mi supervisión directa.
        \end{itemize}
    \end{block}
    \begin{alertblock}{Lo que no proviene de IA}
        Todas las cifras reportadas provienen de ejecución real del
        código sobre los datos descritos. El diseño del framework,
        la formulación matemática del motor de perturbaciones, la
        interpretación de resultados y las conclusiones son
        responsabilidad exclusiva de mi persona.
    \end{alertblock}
\end{frame}

% -----------------------------------------------------------------------------
\begin{frame}{Referencias Principales}
    \tiny
    \begin{thebibliography}{99}
        \bibitem{arsenault2025} Arsenault et al.\ (2025). Revisión
        sistemática de XAI en series temporales financieras.
        \bibitem{tonekaboni2020} Tonekaboni et al.\ (2020). FIT:
        Feature Importance in Time.
        \bibitem{leung2023winit} Leung et al.\ (2023). WinIT.
        \bibitem{crabbe2021} Crabbé y Van Der Schaar (2021). DynaMask.
        \bibitem{enguehard2023} Enguehard (2023). Learning
        Perturbations to Explain Time Series Predictions (ExtrMask).
        \bibitem{liu2024contralsp} Liu et al.\ (2024). ContraLSP.
        \bibitem{hooker2021} Hooker, Mentch y Zhou (2021). Unrestricted
        Permutation Forces Extrapolation.
        \bibitem{bento2021} Bento et al.\ (2021). TimeSHAP.
        \bibitem{raykar`2023} Raykar et al.\ (2023). TsSHAP.
        \bibitem{roshinta2026} Roshinta et al.\ (2026). Ventanas de
        perturbación dinámicas.
        \bibitem{gsshap2026} Kim y Park (2026). GS-SHAP.
        \bibitem{schlegel2019} Schlegel et al.\ (2019). Evaluación de
        métodos XAI en series temporales.
        \bibitem{huang2025} Huang et al.\ (2025). ShapeX.
        \bibitem{yadav2025} Yadav y Subbian (2025).
        \bibitem{adebayo2018} Adebayo et al.\ (2018). Sanity checks
        for saliency maps.
        \bibitem{samek2017aopc} Samek et al.\ (2017). AOPC.
        \bibitem{captum2020} Captum (2020). FeatureAblation.
    \end{thebibliography}
\end{frame}

\end{document}
